% problem-000237 8 有机化合物基本性质
\section*{8. 有机化合物基本性质}
\textit{下列为分问。}

\subsection*{8-1}
下列化合物中酸性最强的是: ( )

（图：）

(a)

（图：）

(b)

（图：）

(c)

（图：）
(d)

\subsection*{8-2}
用化学方法鉴别下列化合物，需要选用的试剂有：（）

（图：）

（图：）

（图：）

(a) HCl 溶液；(b) 溴水；(c) 对甲苯磺酰氯；(d) NaOH 水溶液；(e) $CH_{3}I$

\subsection*{8-3}
下图示出了四个有机化合物指定碳氢键在二甲亚砜(DMSO)中的 $\mathrm{pK}_{\mathrm{a}}$ 值。选择合适的答案解释以下问题:

（图：）

\subsection*{8-3}
-1 氮原子的电负性大于硫，而化合物 3 的酸性却显著大于 1，其原因在于：（）

(a) 化合物 3 形成的负离子稳定；(b) 碳负离子孤对电子对 C-S 键反键轨道的超共轭效应比 C-N 键的强；
(c) 噬吩环的芳香性弱； (d) 以上答案都不对。

8-3-2 氧原子的电负性大于硫，而化合物 3 的酸性却大于 4，其原因在于：（）
(a) 化合物 3 形成的负离子稳定； (b) 噻吩中 C2 位的 C-H 键中 s 成分高；
(c) 噻吩环的吸电子作用更强； (d) 以上答案都不对。

8-3-3 你认为化合物 5 的 $pK_{a}$ 值大概是: ( )
(a) 介于 2 和 3 之间; (b) 介于 3 和 4 之间; (c) 小于 3 的 $pK_{a}$ 值; (d) 以上答案都不对。

8- 4 有人想用 2-溴噻吩和乙炔钠进行反应，结果意外检测到了反应的副产物。你认为副产物可能是：（）
(a) 2,5-二溴噻吩； (b) 四氢噻吩； (c) 噻吩； (d) 以上答案都不对。

\subsection*{8-5}
聚脲是一类高性能聚合物树脂材料，广泛应用于防水防腐涂料、胶黏剂和浇铸材料等领域；而聚硫脲具有优异的自修复性能、高折光指数、高介电常数、强金属配位能力等优势。这两种材料均可以通过链间的氢键作用增强其性能。以下说法正确的有：（）

（图：）

（图：）

(a) 聚合物 A 分子间氢键更强； (b) 聚合物 B 分子间氢键更强； (c) 无法比较。

\subsection*{8-6}
某化合物酮 A 的化学式为 $C_{7}H_{14}O$ ，它在 NaOH 溶液中与 $I_{2}$ 反应后，生成 B 和一黄色沉淀。B 经酸化处理提纯后，再用 P 和 $Br_{2}$ 处理，再经 NaOH 处理生成 C，C 用催化量的高锰酸钾和 $NaIO_{4}$ 溶液处理后得到分子量为 74 的一元酸 D 和另一个一元酸 E。画出化合物 A 的结构式。
